\documentclass[11pt,letterpaper]{article}

\usepackage[margin=1in]{geometry}
\usepackage{amsmath, amssymb}
\usepackage{tikz}
\usepackage{enumitem}
\setlength{\parindent}{0pt}
\setlength{\parskip}{1ex}

\usepackage{hyperref}
\usepackage{listings}
\usepackage{xcolor}
\usepackage{longtable}

\title{Enterprise Building Directory Kiosk Deployment Manual\\Debian 13 (Trixie)}
\author{}
\date{\today}

\begin{document}
\maketitle

\section{Purpose}

This document defines the enterprise deployment, hardware configuration, and long-term maintenance procedures for the Debian 13 building directory kiosk fleet.

Scope includes:

\begin{itemize}
\item Initial hardware provisioning
\item BIOS hardening
\item Debian 13 installation
\item Read-only resilient kiosk configuration
\item Development vs production server placement
\item Multi-kiosk fleet deployment
\item Recovery and maintenance procedures
\end{itemize}

Audience: IT administrators and facilities technology staff.

\subsection{Current-vs-Future Scope}

\begin{itemize}
\item This guide's operational commands target the \textbf{current implemented model}.
\item Any future architecture items are explicitly marked as planned/non-implemented.
\item If a future section conflicts with current scripts, follow current scripts.
\end{itemize}

\section{Hardware Platform}

\subsection{Primary Platform: Qotom Q305P}

Model: Qotom Q305P Mini Desktop PC  
CPU: Intel 3205U  
RAM: 4GB  
Storage: 32GB SSD  
Cooling: Fanless  
Ports: HDMI, VGA, multiple USB, dual NIC, COM  

\subsection{Rationale for Selection}

\begin{itemize}
\item Fanless for 24/7 public display reliability
\item Low power draw (<15W typical)
\item No moving parts
\item Dual NIC (future VLAN segmentation possible)
\item Compact mounting behind display
\item Proven Linux compatibility
\end{itemize}

\subsection{Minimum Hardware Requirements}

\begin{longtable}{p{3in} p{2in}}
CPU & Dual core x86\_64 \\ 
RAM & 4GB minimum (8GB recommended future) \\ 
Storage & 32GB SSD minimum \\ 
GPU & Integrated Intel supported by Debian \\ 
Network & Gigabit Ethernet \\ 
\end{longtable}

\subsection{Recommended Future Hardware}

For next refresh cycle:

\begin{itemize}
\item Intel N100 or N305 fanless mini PC
\item 8GB RAM
\item 128GB SSD
\item Dual HDMI
\item TPM optional
\item Passive cooling only
\end{itemize}

\section{BIOS Configuration (AMI v5.010)}

Enter BIOS: press DEL at boot.

\subsection{Boot}

\begin{longtable}{p{3in} p{2in}}
Boot Mode & UEFI \\
Secure Boot & Disabled \\
Fast Boot & Disabled \\
Boot from USB & Disabled (production) \\
PXE Boot & Disabled \\
Boot order & Internal SSD only \\
\end{longtable}

\subsection{Power}

\begin{longtable}{p{3in} p{2in}}
Restore after AC loss & Power On \\
ErP & Disabled \\
Wake timers & Disabled \\
Sleep states & Disabled \\
\end{longtable}

\subsection{Security}

\begin{longtable}{p{3in} p{2in}}
Admin password & Set and documented \\
External boot & Disabled \\
TPM & Disabled unless required \\
\end{longtable}

\section{Debian 13 Installation}

Install Debian 13 minimal.

Select only:
\begin{itemize}
\item Standard system utilities
\item SSH server
\end{itemize}

Create user:
\begin{lstlisting}
kiosk
\end{lstlisting}

\section{Base Packages}

\begin{lstlisting}[language=bash]
sudo apt update
sudo apt install \
xorg xfce4 \
chromium cage wlr-randr \
overlayroot curl vim
\end{lstlisting}

\textbf{Note:} Canonical package installation and system setup are defined by:
\begin{itemize}
\item \texttt{building-directory-install/install.sh}
\item \texttt{building-directory-install/deploy.sh}
\end{itemize}

\section{Read-Only Resilience Configuration}

\subsection{overlayroot}

\begin{lstlisting}[language=bash]
sudo vi /etc/overlayroot.conf
\end{lstlisting}

Set:
\begin{lstlisting}
overlayroot="tmpfs:swap=1,recurse=0"
\end{lstlisting}

The \texttt{recurse=0} flag is required so the persistent \texttt{/data}
partition remains outside the overlay and retains database writes across reboot.

Reboot and verify:
\begin{lstlisting}[language=bash]
mount | grep overlayroot
\end{lstlisting}

\section{Maintenance Boot Mode}

Add GRUB entry:

\begin{lstlisting}
menuentry "MAINTENANCE MODE (RW)" {
 linux /vmlinuz root=/dev/sda2 ro overlayroot=disabled
 initrd /initrd.img
}
\end{lstlisting}

\begin{lstlisting}[language=bash]
sudo update-grub
\end{lstlisting}

\section{Autologin and Kiosk Launch}

Disable GUI manager:

\begin{lstlisting}[language=bash]
sudo systemctl disable lightdm --now
sudo systemctl set-default multi-user.target
\end{lstlisting}

TTY autologin:

\begin{lstlisting}
[Service]
ExecStart=
ExecStart=-/sbin/agetty --autologin kiosk --noclear %I $TERM
Type=simple
\end{lstlisting}

\section{Chromium Kiosk Launch}

\begin{lstlisting}[language=bash]
#!/bin/bash
SERVER_URL="http://localhost"

exec cage -d -- sh -c '
exec chromium \
--ozone-platform=wayland \
--kiosk \
--no-first-run \
--disable-session-crashed-bubble \
--disable-infobars \
--disable-translate \
'"$SERVER_URL"'
\end{lstlisting}

\section{Keyboard Maintenance Mode}

When USB keyboard inserted:

\begin{itemize}
\item kiosk exits
\item XFCE launches
\item Admin performs maintenance
\item Logout returns to kiosk
\end{itemize}

All XFCE state redirected to /tmp.

\section{Server Placement Strategy}

\subsection{Development Mode}

Server on development workstation:
\begin{lstlisting}
http://192.168.1.131/
http://192.168.1.131/admin
\end{lstlisting}

Local API health checks may still use localhost port 3000:
\begin{lstlisting}[language=bash]
curl -fsS http://127.0.0.1:3000/api/data-version
\end{lstlisting}

\subsection{Production Mode}

Primary server hosted on the production kiosk node:
\begin{lstlisting}
http://192.168.1.80/
http://192.168.1.80/admin
\end{lstlisting}

Standby server:
\begin{lstlisting}
http://192.168.1.81/
http://192.168.1.81/admin
\end{lstlisting}

Kiosks are deployed with a primary URL of \texttt{http://192.168.1.80} and a
standby URL of \texttt{http://192.168.1.81}. On startup they try the primary
first and fall back to the standby if the primary is unreachable. They do not
automatically fail back to the primary until the kiosk session is restarted.

\subsection{Do Not Host Server On All Displays}

One server only.

Reason:
\begin{itemize}
\item Prevent data divergence
\item Simplify updates
\item Easier backup
\end{itemize}

\section{Current Runtime and Operations Model (Implemented)}

\subsection{Server Runtime}

\begin{itemize}
\item Current systemd service name: \texttt{directory-server}
\item Server path: \texttt{/home/kiosk/building-directory/server}
\item App entrypoint: \texttt{server.js}
\item nginx is required for standard port-80 access to \texttt{/} and \texttt{/admin}
\item Admin logout behavior: selecting logout in \texttt{/admin} redirects to base URL \texttt{/}
\item Environment is managed by installer and service overrides (if needed)
\end{itemize}

\subsection{Nginx Prerequisite Note}

When static files are served from \texttt{/home/kiosk/building-directory/...}, nginx must be able to traverse \texttt{/home/kiosk}:
\begin{lstlisting}[language=bash]
sudo chmod 711 /home/kiosk
\end{lstlisting}

This preserves home privacy while allowing nginx path traversal.

\subsection{Kiosk Client Runtime}

\begin{itemize}
\item Kiosk loop launches \texttt{cage + chromium} from tty1 autologin
\item USB keyboard insertion exits kiosk and starts XFCE fallback
\item Kiosk system scripts are updated via server-side deploy workflow
\end{itemize}

\subsection{Current Day-2 Deployment Flows}

\begin{itemize}
\item \textbf{Dev machine == deployed server:} use \texttt{tools/deploy-local.sh}
\item \textbf{Deploy server code to target machine:} use \texttt{building-directory-install/deploy.sh}
\item \textbf{Deploy kiosk system scripts to display clients:}
  use Admin Deploy tab (\texttt{POST /api/kiosks/:id/deploy}) or \texttt{server/kiosk-deploy.sh}
\end{itemize}

\subsection{Current Health/Verification Commands}

\begin{lstlisting}[language=bash]
sudo systemctl status directory-server
sudo systemctl restart directory-server
curl -fsS http://127.0.0.1:3000/api/data-version
\end{lstlisting}

\section{Firewall Rules}

Server system:

\begin{lstlisting}[language=bash]
sudo ufw allow 22
sudo ufw allow from 192.168.1.0/24 to any port 80
# Optional (only if direct Node access is intentionally required):
# sudo ufw allow from 192.168.1.0/24 to any port 3000
sudo ufw enable
\end{lstlisting}

Best practice: keep port 3000 internal and serve kiosk/admin traffic through nginx on port 80.

\section{Fleet Deployment Plan}

Clone master image after validation.

Steps:

\begin{enumerate}
\item Install first kiosk fully
\item Validate power-loss recovery
\item Validate overlayroot
\item Clone disk to other units
\item Assign static IPs
\item Update hostname
\end{enumerate}

\section{Operational Validation}

\begin{enumerate}
\item Cold boot → kiosk auto starts
\item Power loss recovery works
\item Chromium launches full screen
\item Server reachable
\item Keyboard triggers admin mode
\item Logout returns to kiosk
\end{enumerate}

\section{Failure Recovery}

\subsection{If kiosk frozen}

Power cycle.

\subsection{If corrupted config}

Boot maintenance mode.

\subsection{If server unreachable}

Check:
\begin{itemize}
\item network link
\item firewall
\item node server running
\end{itemize}

\section{Change Control and Deployment Best Practices}

\subsection{Source of Truth and Promotion}

\begin{itemize}
\item Develop and review in \texttt{Public-Kiosk}.
\item Promote to deployed environment only via documented deploy scripts.
\item Avoid manual file edits on production nodes outside maintenance mode.
\end{itemize}

\subsection{Pre-Deployment Checklist}

\begin{enumerate}
\item Confirm target environment (dev vs production)
\item Confirm current git commit and pending local changes
\item Run dry-run where supported:
\begin{lstlisting}[language=bash]
tools/deploy-local.sh --dry-run
\end{lstlisting}
\item Verify required service account access (SSH keys, sudo rules)
\item Confirm backup/restore path is available
\end{enumerate}

\subsection{Post-Deployment Checklist}

\begin{enumerate}
\item Service active:
\begin{lstlisting}[language=bash]
systemctl is-active directory-server
\end{lstlisting}
\item API health:
\begin{lstlisting}[language=bash]
curl -fsS http://127.0.0.1:3000/api/data-version
\end{lstlisting}
\item UI checks: kiosk main screen, company/individual search, building info
\item Admin checks: login, CSV upload/download, deploy tab
\item Kiosk checks: keyboard maintenance flow and return to kiosk
\end{enumerate}

\subsection{Rollback (Current Model)}

\begin{enumerate}
\item Revert source to known-good commit/tag
\item Re-run deploy script (\texttt{tools/deploy-local.sh} or target deploy workflow)
\item Restart service:
\begin{lstlisting}[language=bash]
sudo systemctl restart directory-server
\end{lstlisting}
\item Re-run health and UI validation checks
\end{enumerate}

\subsection{Overlayroot/Maintenance Mode Note}

Kiosk script deploys are designed for normal overlayroot runtime.  
If a kiosk is booted with \texttt{overlayroot=disabled} (maintenance mode), standard overlayroot-chroot deploy helpers may fail and direct filesystem updates may be required.

\subsection{Backup and Restore Practice}

\begin{itemize}
\item Capture regular backups before major deploys.
\item Preferred admin backup download is SQL text: \texttt{GET /api/backup.txt} (\texttt{directory-backup.txt}).
\item Additional SQL dump endpoint is available: \texttt{GET /api/backup.sql}.
\item Restore endpoint \texttt{POST /api/restore} accepts \texttt{.txt}, \texttt{.sql}, \texttt{.sqlite}, and \texttt{.db}.
\item Verify restore procedure periodically.
\item Keep at least one known-good backup off-device.
\end{itemize}


\section{Current Update System (Admin Workstation and Server Deploy)}

This section documents the update process that is currently implemented in the repository and deployment environment.

\subsection{Design Goals}

\begin{itemize}
\item Updates are initiated from development/admin machine workflows
\item One server node remains the data source for all displays
\item Display nodes auto-consume client UI updates from the server
\item Kiosk machine script updates are pushed explicitly via SSH deploy
\end{itemize}

\subsection{Implemented Components}

\begin{itemize}
\item Local deploy script: \texttt{tools/deploy-local.sh}
\item Remote SSH deploy script: \texttt{tools/deploy-ssh.sh}
\item Remote/full deploy script: \texttt{building-directory-install/deploy.sh}
\item Kiosk script deploy script: \texttt{server/kiosk-deploy.sh}
\item Admin UI deploy endpoints:
\begin{itemize}
\item \texttt{GET /api/kiosks}
\item \texttt{GET /api/kiosks/deploy-pubkey}
\item \texttt{POST /api/kiosks/:id/deploy}
\end{itemize}
\end{itemize}

\subsection{Typical Current Workflow}

\begin{enumerate}
\item Make and verify changes in \texttt{Public-Kiosk}
\item For remote server nodes, preview and deploy:
\begin{lstlisting}[language=bash]
tools/deploy-ssh.sh --dry-run --host kiosk@192.168.1.80
tools/deploy-ssh.sh --host kiosk@192.168.1.80
# Optional data promotion:
tools/deploy-ssh.sh --host kiosk@192.168.1.80 \
  --with-db --db-source /home/security/building-directory/server/directory.db
\end{lstlisting}
\item Ensure \texttt{rsync} is installed on both source and target hosts:
\begin{lstlisting}[language=bash]
sudo apt-get update && sudo apt-get install -y rsync
\end{lstlisting}
\item Preview deploy impact:
\begin{lstlisting}[language=bash]
tools/deploy-local.sh --dry-run
\end{lstlisting}
\item Deploy server changes:
\begin{lstlisting}[language=bash]
tools/deploy-local.sh
\end{lstlisting}
\item Restart service if prompted:
\begin{lstlisting}[language=bash]
sudo systemctl restart directory-server
\end{lstlisting}
\item Deploy kiosk script updates from Admin Deploy tab (or \texttt{kiosk-deploy.sh})
\item Validate health:
\begin{lstlisting}[language=bash]
curl -fsS http://127.0.0.1:3000/api/data-version
systemctl is-active directory-server
journalctl -u directory-server -n 80 --no-pager
\end{lstlisting}
\end{enumerate}

\subsection{Fleet Consideration}

Display nodes generally do not require direct content deployment for kiosk HTML/CSS/JS changes; they load client assets from the server URL.

System-script changes (\texttt{start-kiosk.sh}, \texttt{.bash\_profile}, udev rules, power-button policy) do require explicit kiosk deploy.

\section{Testing Procedures}

\subsection{Purpose}

This section defines repeatable test procedures with explicit pass/fail criteria for development, deployment, and production readiness.

\subsection{Test Environments}

\begin{longtable}{p{2.2in} p{3.1in}}
Development Server & 192.168.1.131 \\
Production Server & 192.168.1.80 \\
Kiosk Clients & 192.168.1.80, 192.168.1.81, 192.168.1.82 \\
Admin URL & \texttt{http://<server-ip>/admin} \\
Kiosk URL & \texttt{http://<server-ip>/} \\
\end{longtable}

\subsection{Entry Criteria}

\begin{itemize}
\item Target host is reachable via SSH.
\item \texttt{directory-server} service is installed.
\item Correct branch/commit is deployed.
\item Backup is taken before production deployment.
\end{itemize}

\subsection{Core Health Checks}

\begin{lstlisting}[language=bash]
systemctl is-active directory-server
curl -fsS http://127.0.0.1:3000/api/data-version
curl -fsS http://127.0.0.1:3000/api/auth/me
journalctl -u directory-server -n 80 --no-pager
\end{lstlisting}

\textbf{Pass criteria:}
\begin{itemize}
\item Service state is \texttt{active}.
\item API endpoints return HTTP 200.
\item No crash loop or fatal errors in recent logs.
\end{itemize}

\subsection{Functional Test Matrix}

\begin{longtable}{p{0.6in} p{2.0in} p{1.6in} p{1.7in}}
ID & Test & Method & Pass Criteria \\
\hline
T1 & Kiosk auto-start after reboot & Reboot client & Cage + Chromium launch fullscreen within 60s \\
T2 & Main UI render & Visual check & Welcome banner, buttons, background render correctly \\
T3 & Company search & Use on-screen keyboard & Matching results shown; no UI error \\
T4 & Individual search & Use on-screen keyboard & Matching results shown; building number shown \\
T5 & Building info screen & Open from main menu & Static content renders correctly \\
T6 & Keyboard maintenance mode & Insert USB keyboard & Kiosk exits, XFCE launches \\
T7 & Return to kiosk & Logout XFCE & Kiosk relaunches automatically \\
T8 & Admin login success & Login with admin password & Authenticated session established \\
T9 & Admin login failure/lockout & Repeated wrong password & 401 then lockout/rate-limit response \\
T10 & Admin logout redirect & Click logout in admin UI & Session cleared; browser redirects to \texttt{/} \\
T11 & Text backup export & Download Text Backup & \texttt{directory-backup.txt} downloads successfully \\
T12 & SQL text restore & Upload valid \texttt{.txt} SQL dump & Restore succeeds; expected DB counts returned \\
T13 & Background image upload & Upload image in admin Appearance tab & Upload succeeds and new image appears in gallery \\
T14 & Unauthorized admin access & Request protected endpoint without auth & HTTP 401 Unauthorized \\
T15 & Companies CSV export & Download CSV & File downloads; header row present even with 0 rows \\
T16 & Companies CSV import & Upload valid CSV & Success status and expected row count in DB \\
T17 & Individuals CSV import & Upload valid CSV & Success status and expected row count in DB \\
T18 & Unauthorized CSV access & Request CSV endpoint without auth & HTTP 401 Unauthorized \\
T19 & Data version change & Modify company/individual & \texttt{/api/data-version} increments \\
T20 & Deploy kiosk scripts & Admin deploy tab or script & Deploy reports success; target script updated \\
T21 & Power button long press & Hold power button on client & System powers off cleanly \\
\end{longtable}

\subsection{Resilience and Negative Tests}

\begin{longtable}{p{0.6in} p{2.0in} p{3.3in}}
ID & Scenario & Pass Criteria \\
\hline
R1 & Server restart during kiosk operation & Kiosk reconnects and resumes normal operation \\
R2 & Network interruption (temporary) & UI recovers after network returns; no permanent lock-up \\
R3 & Overlay maintenance mode deploy attempt & Procedure detects mode and uses correct deploy path \\
R4 & Empty companies export & CSV download still returns valid header row \\
R5 & Invalid CSV upload & Clear error message; no partial/invalid DB write \\
R6 & Invalid SQL/text restore upload & Restore rejected with clear error; no partial DB replacement \\
\end{longtable}

\subsection{Deployment Validation (Per Release)}

\begin{enumerate}
\item Run core health checks on target server.
\item Execute T1--T8 minimum.
\item Execute T10--T14 for auth/backup/admin workflow changes.
\item Execute T15--T18 for CSV workflow changes.
\item Execute T20 for any kiosk script changes.
\item Record results in deployment log.
\end{enumerate}

\subsection{Rollback Trigger Criteria}

Rollback is required if any of the following occur:
\begin{itemize}
\item \texttt{directory-server} fails to remain active.
\item Kiosk UI fails to render or is non-interactive.
\item Admin auth, backup/restore, or CSV workflows fail.
\item Data integrity checks fail (missing/incorrect records after import).
\end{itemize}

\subsection{Test Evidence Template}

\begin{longtable}{p{1.1in} p{1.0in} p{0.7in} p{1.3in} p{1.7in}}
Date/Time & Tester & Test ID & Result (Pass/Fail) & Notes / Evidence \\
\hline
 &  &  &  &  \\
 &  &  &  &  \\
 &  &  &  &  \\
\end{longtable}


\section{Production Signoff}

System approved when:

\begin{itemize}
\item BIOS locked
\item Read-only root active
\item Auto-start verified
\item Network stable
\item 72-hour continuous run test passed
\end{itemize}

\section{Future Expansion}

Planned enhancements:

\begin{itemize}
\item Central management node
\item Health heartbeat monitoring
\item Secondary backup server
\item VLAN isolation
\end{itemize}

\section{Future Architecture (Planned, Not Yet Implemented)}

The following model is a proposed future-state architecture and should not be treated as current operations until scripts and services are migrated.

\begin{itemize}
\item \texttt{/srv/building-directory/releases/<timestamp>} layout
\item \texttt{current} symlink for atomic release switching
\item Separate persistent shared directories (\texttt{db}, \texttt{uploads}, \texttt{logs})
\item Optional dedicated service (example: \texttt{building-directory.service})
\item rsync release push with symlink rollback
\end{itemize}

When this migration is executed, update this guide and the deploy scripts together in the same release.

\end{document}
